\chapter{Making the comparison defensible}\label{ch:comparison}
\section{Explain why these jurisdictions belong together}
State the purpose of the comparison before choosing the jurisdictions.
Describe the feature being compared and the institutional conditions that make
the comparison informative. Similar terminology may conceal different powers,
procedures or remedies. Different terminology may describe arrangements that
answer a comparable problem.

Comparative research also depends on access to language and contextual
knowledge. Explain how these considerations affected selection and interpretation.
\lawcite[3--4]{vanhoecke2015}

\section{Compare a common question across the materials}
Use the same analytical question in each jurisdiction while preserving legally
significant differences. Record who exercises the relevant power, what triggers
it, who may contest it, which forum hears the dispute and what consequences
follow. The accompanying comparative matrix provides editable fields for these
questions.

\LawPrompt{Populate the matrix from verified sources. Mark an unanswered
question as unresolved instead of treating the absence of a source as proof
that no rule or practice exists.}

\section{Explain what the comparison establishes}
Separate observed differences from explanations for them. A doctrinal comparison
can establish different legal conditions without demonstrating why institutions
behave differently in practice. Where an explanatory claim depends on interviews,
case selection or quantitative evidence, supply the appropriate method and
limits. Acknowledge alternatives that the available evidence cannot resolve.
